% ============================================================
% §4 模型假设与合理性论证 (Stage 04 产物)
% 竞赛: 电工杯 B题 (2026)
% ============================================================

\section{模型假设与合理性论证}

建模本质上是对现实世界的受控简化——假设的质量直接决定模型的解释力边界。本文在充分尊重题目背景与附件数据结构的前提下，提炼以下四条核心假设，每条均附类型分类与支撑依据，避免反模式B1（假设无支撑）与B2（假设过多导致模型脆弱）。

\begin{table}[htbp]
\centering
\caption{核心模型假设清单}
\label{tab:assumptions}
\begin{tabular}{p{0.8cm}p{5.5cm}p{2.2cm}p{4.5cm}}
\toprule
\textbf{编号} & \textbf{假设内容} & \textbf{类型} & \textbf{支撑依据} \\
\midrule
A1 & 各小区老年人口的马尔可夫状态转移服从独立同质的年际递推规律，即自理$\to$半失能$\to$失能的退化过程在所有10个小区中遵循相同的转移概率矩阵$\mathbf{A}$，且各小区之间不存在老年人口的跨区迁移 & 模型隐含 & 题目明确给出统一转移概率（附件1），且10个小区隶属于同一成熟街道，社区间的人口流动性在微观尺度上可忽略；独立同质假设是马尔可夫链的标准前提 \\
\midrule
A2 & 五年预测窗口内，当地工资水平、物价指数与政府养老补贴政策保持平稳，不发生结构性突变 & 环境假设 & 题目明确要求“当地工资、物价与补贴政策保持稳定”；该假设将附件2--3的静态成本参数（基准价、建设成本、补贴标准）的有效期从基年合理外推至第5年末，避免引入不可观测的政策冲击 \\
\midrule
A3 & 老人选择满意度的行为完全由附件5的三因素加权评分函数$S=0.2S_1+0.3S_2+0.5S_3$刻画，不存在未观测的偏好异质性（如品牌忠诚度、邻里口碑等非结构化因素） & 模型隐含 & 附件5作为题目唯一给定的满意度评价框架，其评分权重的设定蕴含着出题方对养老服务质量核心维度的价值判断；将老人行为锚定于该函数，保证了模型的可验证性与可复现性 \\
\midrule
A4 & 服务站的日最大服务能力$Q_k^{\max}$为硬性上限（不可超载），且日均固定管理成本$O_k$不随实际服务人次的波动而变化（即成本函数为分段常数：建站前=0，建站后=固定值） & 数据假设 & 附件3对三种规模的“日最大服务人次”采用确定性上限表述（无概率分布），符合工程设施设计中“额定容量”的行业惯例；固定管理成本对应人员工资、场地租金等刚性支出 \\
\bottomrule
\end{tabular}
\end{table}

\textbf{假设A1的深度论证}：将10个小区的老年人口演化建模为独立同质的马尔可夫链，是全文数理基座的逻辑起点。其数学合理性体现在三个层面：(1) 退化方向性——题设明确“失能不能恢复为自理或半失能”“半失能不能恢复为自理”，这恰好对应马尔可夫链的“部分吸收”结构（矩阵$\mathbf{A}$的零元位置$[1,2]=[2,0]=[2,2]=0$严格编码了不可逆性）；(2) 增量开放性——每年7\%的新增自理老人从“外部环境”注入系统，使转移矩阵的列和大于1（$\sum_j a_{ji} > 1$），形成一种“增生型马尔可夫过程”，区别于标准闭合马尔可夫链的人口递减趋势；(3) 空间独立性——10个小区间无迁移项，使得10组递推可完全并行计算（矩阵相同，初始向量不同），大幅降低了计算复杂度而几乎不损失建模精度——因为社区养老服务站的“嵌入式”定位本就意味着服务的社区内闭环特征。

\textbf{假设A3的补充说明}：附件5的综合满意度函数$S = 0.2S_1 + 0.3S_2 + 0.5S_3$中，价格满意度$S_3$的权重高达0.5——这反映出题设对“服务可负担性”的高度重视。在Q3的定价优化中，我们将显式利用这一权重结构：通过降低服务价格提升$S_3$，但需在利润率$\le 8\%$的监管约束下寻求帕累托平衡点。这种“权重驱动的目标导向”正是将评分规则从描述性工具转化为规范性优化目标的关键一步。
